% ==============================================================================
% Four-Winding Transformer — Circuit Model and Equations
% Author: Davide Bagnara
% ==============================================================================
\documentclass[11pt,a4paper]{article}

% --- Packages -----------------------------------------------------------------
\usepackage[T1]{fontenc}
\usepackage[utf8]{inputenc}
\usepackage[english]{babel}
\usepackage{amsmath,amssymb,bm}
\usepackage{booktabs}
\usepackage{array}
\usepackage{geometry}
\usepackage{hyperref}
\usepackage{parskip}
%\usepackage[activate={true,nocompatibility},final,tracking=false,kerning=true,spacing=false]{microtype}
\usepackage{xcolor}

\geometry{
  a4paper,
  top=25mm, bottom=25mm,
  left=25mm, right=25mm
}

\hypersetup{
  colorlinks = true,
  linkcolor  = blue!60!black,
  urlcolor   = blue!60!black
}

% --- Custom commands ----------------------------------------------------------
\newcommand{\mat}[1]{\mathbf{#1}}
\newcommand{\Rel}{\mathcal{R}}          % reluctance
\newcommand{\FMM}{\mathcal{F}}          % magnetomotive force
\newcommand{\der}[2]{\dfrac{d#1}{d#2}}

% ==============================================================================
\begin{document}

\title{\textbf{Four-Winding Transformer}\\[4pt]
       \large Circuit Model and Governing Equations}
\author{Davide Bagnara}
\date{\today}
\maketitle
\tableofcontents
\vspace{6pt}
\hrule
\vspace{10pt}

% ------------------------------------------------------------------------------
\section{Introduction}

This document derives the governing equations of a four-winding transformer
sharing a common magnetic core.
The four windings are grouped as two \emph{primary} (indices $1$, $2$) and two
\emph{secondary} (indices $3$, $4$) windings.
All windings are magnetically coupled through the same flux path in the core.

Notation used throughout:
\begin{itemize}
  \item $N_k$ — number of turns of winding $k$
  \item $v_k$, $i_k$, $\lambda_k$ — terminal voltage, current, and flux linkage
        of winding $k$
  \item $R_k$ — winding resistance
  \item $L_{kk}$ — self-inductance of winding $k$
  \item $M_{jk} = M_{kj}$ — mutual inductance between windings $j$ and $k$
  \item $\Rel$ — core reluctance
  \item $\Phi_m$ — common magnetising flux
\end{itemize}

% ------------------------------------------------------------------------------
\section{Flux Linkage Matrix}

Each winding concatenates the total core flux. The flux linkage of winding $k$
is:
\begin{equation}
  \lambda_k = L_{kk}\,i_k + \sum_{j \neq k} M_{kj}\,i_j,
  \qquad k = 1,2,3,4.
  \label{eq:flux_scalar}
\end{equation}

In compact matrix form:
\begin{equation}
  \bm{\lambda} = \mat{L}\,\mathbf{i},
  \label{eq:flux_matrix}
\end{equation}

where
\begin{equation}
  \begin{bmatrix}
    \lambda_1 \\ \lambda_2 \\ \lambda_3 \\ \lambda_4
  \end{bmatrix}
  =
  \begin{bmatrix}
    L_{11} & M_{12} & M_{13} & M_{14} \\
    M_{12} & L_{22} & M_{23} & M_{24} \\
    M_{13} & M_{23} & L_{33} & M_{34} \\
    M_{14} & M_{24} & M_{34} & L_{44}
  \end{bmatrix}
  \begin{bmatrix}
    i_1 \\ i_2 \\ i_3 \\ i_4
  \end{bmatrix}.
  \label{eq:flux_full}
\end{equation}

The inductance matrix $\mat{L}$ is \emph{symmetric} and \emph{positive definite}.

\subsection{Ideal Core — Perfect Coupling}

For a lossless, infinite-permeability core (no leakage), all mutual inductances
satisfy:
\begin{equation}
  M_{jk} = \frac{N_j N_k}{\Rel} = \sqrt{L_{jj}\,L_{kk}},
  \label{eq:mutual_ideal}
\end{equation}

and the coupling coefficient between any pair equals unity:
\begin{equation}
  k_{jm} \triangleq \frac{M_{jm}}{\sqrt{L_{jj}\,L_{mm}}} = 1.
\end{equation}

% ------------------------------------------------------------------------------
\section{Voltage Equations (KVL)}

Applying Faraday's law and Kirchhoff's voltage law to each winding:
\begin{equation}
  \boxed{v_k = R_k\,i_k + \der{\lambda_k}{t}
         = R_k\,i_k + L_{kk}\der{i_k}{t}
           + \sum_{j \neq k} M_{kj}\der{i_j}{t}},
  \qquad k=1,2,3,4.
  \label{eq:voltage_scalar}
\end{equation}

In vector–matrix notation:
\begin{equation}
  \mathbf{v} = \mat{R}\,\mathbf{i} + \mat{L}\,\der{\mathbf{i}}{t},
  \label{eq:voltage_matrix}
\end{equation}

where $\mat{R} = \operatorname{diag}(R_1, R_2, R_3, R_4)$.

% ------------------------------------------------------------------------------
\section{T-Equivalent Model with Leakage and Magnetising Inductances}

\subsection{Inductance Decomposition}

Each self-inductance is decomposed into a \emph{leakage} term $L_{\sigma k}$
and a \emph{magnetising} term referred to winding $k$:
\begin{equation}
  L_{kk} = L_{\sigma k} + \frac{N_k^2}{\Rel} = L_{\sigma k} + L_{m,k},
  \label{eq:decomp}
\end{equation}

where $L_{m,k} = N_k^2/\Rel$ is the magnetising inductance seen from winding $k$.

\subsection{Magnetising Flux and Back-EMF}

The common magnetising flux is driven by the total magnetomotive force (MMF):
\begin{equation}
  \Phi_m = \frac{\FMM_{tot}}{\Rel}
           = \frac{1}{\Rel}\sum_{k=1}^{4} N_k\,i_k.
  \label{eq:flux_core}
\end{equation}

The back-EMF induced in winding $k$ by the common flux is:
\begin{equation}
  e_k = N_k\,\der{\Phi_m}{t}.
  \label{eq:emf}
\end{equation}

\subsection{Voltage Equations — T Form}

Each winding equation splits into a resistive-leakage drop and a back-EMF:
\begin{equation}
  v_k = R_k\,i_k + L_{\sigma k}\der{i_k}{t} + e_k.
  \label{eq:voltage_T}
\end{equation}

\subsection{Back-EMF Proportionality}

Since $e_k = N_k\,d\Phi_m/dt$, the back-EMFs are proportional to the turns
ratio with respect to any reference winding (here winding 1):
\begin{equation}
  \frac{e_1}{N_1} = \frac{e_2}{N_2} = \frac{e_3}{N_3} = \frac{e_4}{N_4}
  \quad \Longrightarrow \quad
  e_k = \frac{N_k}{N_1}\,e_1.
  \label{eq:emf_ratio}
\end{equation}

% ------------------------------------------------------------------------------
\section{Ampère's Law — MMF Balance}

\subsection{General Form}

The net MMF sustains the core flux against the reluctance:
\begin{equation}
  \sum_{k=1}^{4} N_k\,i_k = \Rel\,\Phi_m.
  \label{eq:ampere_general}
\end{equation}

\subsection{Ideal Core ($\Rel \to 0$, $i_m \approx 0$)}

When the magnetising current is negligible compared to the load currents, the
constraint reduces to:
\begin{equation}
  \boxed{N_1\,i_1 + N_2\,i_2 + N_3\,i_3 + N_4\,i_4 = 0.}
  \label{eq:ampere_ideal}
\end{equation}

This is the generalised Ampère balance for the four-winding ideal transformer.

% ------------------------------------------------------------------------------
\section{Turns Ratios and Port Quantities}

\subsection{Definition}

With winding 1 chosen as the voltage reference:
\begin{equation}
  a_k \triangleq \frac{N_k}{N_1}, \qquad k = 2,3,4.
  \label{eq:turns_ratio}
\end{equation}

\subsection{No-Load Voltages}

Open-circuit terminal voltages obey:
\begin{equation}
  V_k^{(\text{oc})} = a_k\,V_1, \qquad k = 2,3,4.
  \label{eq:voltage_ratio}
\end{equation}

\subsection{Current Constraint (Ideal Core)}

Substituting the turns ratios into \eqref{eq:ampere_ideal}:
\begin{equation}
  i_1 + a_2\,i_2 = -(a_3\,i_3 + a_4\,i_4).
  \label{eq:current_balance}
\end{equation}

The left-hand side represents the combined primary excitation; the right-hand
side is the reflected secondary load.

% ------------------------------------------------------------------------------
\section{Summary of Parameters}

\begin{table}[h!]
  \centering
  \caption{Parameter summary for the four-winding transformer model.}
  \label{tab:parameters}
  \renewcommand{\arraystretch}{1.35}
  \begin{tabular}{@{} l l l @{}}
    \toprule
    \textbf{Symbol} & \textbf{Unit} & \textbf{Description} \\
    \midrule
    $L_{\sigma k}$        & H         & Leakage inductance of winding $k$ \\
    $L_{m,k} = N_k^2/\Rel$ & H       & Magnetising inductance seen from winding $k$ \\
    $M_{jk} = a_j a_k L_{m,1}$ & H   & Mutual inductance between windings $j$ and $k$ \\
    $R_k$                 & $\Omega$  & Winding resistance \\
    $\Rel$                & A/Wb      & Core reluctance (includes air gap if present) \\
    $a_k = N_k/N_1$       & —         & Turns ratio with respect to winding 1 \\
    $e_k$                 & V         & Back-EMF induced by common core flux \\
    $\Phi_m$              & Wb        & Common magnetising flux \\
    \bottomrule
  \end{tabular}
\end{table}

% ------------------------------------------------------------------------------
\section{Remarks on Multi-Port Applications}

\begin{itemize}
  \item The inductance matrix $\mat{L}$ in \eqref{eq:flux_full} is the natural
        starting point for multi-active-bridge (MAB) converter analysis, where
        the leakage inductances $L_{\sigma k}$ act as power-transfer reactances.
  \item For a dual-active-bridge (DAB) extended to four ports, ZVS conditions
        for each switch must be verified independently for each winding current
        at the switching instants.
  \item When two primary windings are fed from different voltage sources
        (e.g.\ two AFE converters), circulating currents between them are
        governed by $(v_1 - v_2)/( L_{\sigma 1} + L_{\sigma 2} + M_{12,\text{leak}})$
        and must be controlled or suppressed by design.
  \item The constraint \eqref{eq:ampere_ideal} allows reduction of the state
        space: given three independent currents, the fourth is determined.
        This is exploited in state-space models for simulation and control
        design.
\end{itemize}

\end{document}
